% Рисунок с доски: обобщённый материал Максвелла — N элементов Максвелла
% параллельно; функция релаксации при N = 2, τ1 << τ2.
\begin{figure}[htbp]
\centering
\begin{tikzpicture}[>={Latex[length=4pt]}, scale=0.95, transform shape]
  \foreach \k in {0, 1.9, 3.8} {
    \begin{scope}[xshift=\k cm]
      \pic[rotate=90] at (0,0.9) {spring};
      \pic[rotate=90] at (0,-0.9) {damper};
      \draw[thick] (0,0) -- (0,0);
      \draw[thick] (0,1.8) -- (0,2.2);
      \draw[thick] (0,-1.8) -- (0,-2.2);
    \end{scope}
  }
  \node at (5.4,0) {$\cdots$};
  \begin{scope}[xshift=7.0cm]
    \pic[rotate=90] at (0,0.9) {spring};
    \pic[rotate=90] at (0,-0.9) {damper};
    \draw[thick] (0,1.8) -- (0,2.2);
    \draw[thick] (0,-1.8) -- (0,-2.2);
  \end{scope}
  \draw[thick] (0,2.2) -- (7.0,2.2);
  \draw[thick] (0,-2.2) -- (7.0,-2.2);
  \draw[thick] (3.5,2.2) -- (3.5,2.7);
  \draw[thick] (3.5,-2.2) -- (3.5,-2.7);

  % график R(t) при N = 2
  \begin{scope}[xshift=9.6cm, yshift=-1.5cm]
    \draw[->] (0,0) -- (5.6,0) node[anchor=north east] {$t$};
    \draw[->] (0,0) -- (0,3.3) node[anchor=south east] {$R$};
    \draw[dashed, gray] (0,2.8) -- (0.25,2.8);
    \node[anchor=east, font=\footnotesize] at (-0.05,2.9) {$2\mu_1 + 2\mu_2$};
    \draw[dashed, gray] (0,1.4) -- (0.8,1.4);
    \node[anchor=east, font=\footnotesize] at (-0.05,1.3) {$2\mu_2$};
    \draw[very thick, blue!60!black, domain=0:5.4, samples=200, smooth]
      plot (\x, {1.4*exp(-\x/0.3) + 1.4*exp(-\x/2.2)});
    \fill (0.8,{1.4*exp(-0.8/0.3) + 1.4*exp(-0.8/2.2)}) circle (1.3pt);
    \node[anchor=north, font=\footnotesize] at (0.8,-0.05) {$\tau_1$};
    \node[anchor=north, font=\footnotesize] at (4.0,-0.05) {$\tau_2$};
  \end{scope}
\end{tikzpicture}
\caption{Обобщённый материал Максвелла: $R(t) = \sum_{k} 2\mu_k e^{-t/\tau_k}$.
При $N=2$ и $\tau_1 \ll \tau_2$ быстрое звено релаксирует почти сразу (излом,
как в обобщённой модели Кельвина–Фойгта), медленное — за время $\tau_2$.}
\label{fig:maxwell-generalized}
\end{figure}
